\section{收缩不是退场（1425--1428）}

永乐帝死后，洪熙朝停止一部分远征和营建，并召还、赦免若干前朝获罪者。诏令能够确认这种优先次序的调整，不能使各地已投入的工役、军费和人事关系立刻消失。洪熙元年五月，\eventanchor{EM-1425-XUANDE-ACCESSION}{仁宗去世，朱瞻基即位为宣德帝}。继位日期无疑；骤逝的原因没有足以裁定的同时代医学证据。新君首先接手的是一套仍由北方供给、交趾驻军和宗室关系拉扯的中枢。

次年，\eventanchor{EM-1426-GAOXU-REBELLION}{汉王朱高煦在乐安起兵}。宣德帝亲征、叛军降伏与其后处置可由朝廷记录确证；叛乱的密谋和处刑细节已被胜者叙事塑形。乐安的道路、城防和附近粮源决定朝廷为何能够迅速把这场冲突限制在一隅，而宗藩问题也由此显出：继承顺利并没有令永乐诸子的资源关系自动归于静止。

\eventanchor{EM-1427-JIAOZHI-WITHDRAWAL}{宣德二年，明军同黎利方面议和，准备撤出交趾}。红河三角洲的驻军需要粮饷，山区的道路与地方抵抗又使官署、卫所和征发不能只靠北京任命维持。\textit{宣宗实录}排出明方的决策顺序，\textit{大越史记全书}则把战争终结置入黎利建国的另一套年序；两者不能彼此抹去。同年撤军承认的是直接统治难以承受的成本，\eventanchor{EM-1428-LE-LOI-RECOGNITION}{翌年调整对黎利政权的朝贡称谓}，则把关系放回册封、使节与边境交涉。所谓“承认”不是一方替另一方裁定合法性，而是双方以不同语言处理一场已无法按旧制维持的战争。

\section{巡抚、船队与幼主（1430--1435）}

\eventanchor{EM-1430-YUQIAN-INSPECTION}{宣德五年，于谦巡抚河南、山西}。黄河—太行之间的灾荒、军政和地方事务无法仅由单一布政司处置，巡抚差遣遂把跨省调度交到一位可直接奏报的官员手中。任命可以核实，奏议中的灾情和成效却是请求行动的文书；仓储、灾异与地方志才可能在有限地点显示粮食和差役如何落下。于谦后来守京的形象，应先放回这一段处理道路、军需和地方官之间关系的经验中理解。

同年，\eventanchor{EM-1430-ZHENGHE-SEVENTH-VOYAGE}{郑和第七次下西洋自南京出航}。刘家港的诏书、仓储、翻译者与船员把官府动员接到季风和港口网络；出航并不等于朝廷能够支配每一段海路。\eventanchor{EM-1431-TIANFEI-STELE}{次年立下的《天妃灵应之记》碑}，是使团为航行、祈祷和奉命所作的自我呈现，能够确认立碑者与年代，不能独证每一次航程的船数和情节。\eventanchor{EM-1433-ZHENGHE-RETURN}{宣德八年舰队返抵中国}，郑和卒年与卒地仍有异说；较可把握的是，长距离使团把南京、沿海港口与印度洋的译者、物资和礼仪联结起来，而不是创造一条由北京单向控制的海上道路。

\eventanchor{EM-1435-XUANDE-DEATH}{宣德十年宣德帝去世，九岁的朱祁镇即位}。太皇太后、内阁、近侍和官僚之间的辅政平衡并非一日形成或一词可以概括为“宦官专权”。谁能向幼主解释边报、军饷和朝贡争端，将在日后的危机中变得格外重要。

\section{滇西的军道与北方的边警（1436--1448）}

\eventanchor{EM-1436-LUCHUAN-CRISIS}{正统元年麓川思任发的扩张引发滇西危机}。北京收到的是边报和调兵请求，云南、贵州的卫所和邻近土司面对的则是山路、向导、粮运与原有联盟。\eventanchor{EM-1438-FIRST-LUCHUAN-CAMPAIGN}{正统三年王骥等首次大举征麓川}，军队深入并不能使土司竞争自动停止；明方战报所称军额、伤亡与“平定”须同地方志、碑刻和异文分开阅读。

\eventanchor{EM-1440-ESEN-RAID}{正统五年也先势力进犯边镇}，使大同、宣府的马匹、仓储和报讯进入中枢账目。北京文书将此归作边防问题，草原行动者则关心赏赐、通道与内部权力；双方的分类不能替代彼此。滇西的军费和北方的警报同时要求资源，正统朝并没有只面对一条边线。

\eventanchor{EM-1441-SECOND-LUCHUAN-CAMPAIGN}{正统六年再度出兵}，\eventanchor{EM-1443-THIRD-LUCHUAN-CAMPAIGN}{八年继续追击、招抚与分置}，显示朝廷在征讨、扶持地方中介和安排行政名目之间反复取舍。远征所耗的不只是军费：云南卫所、驿路和家庭必须提供人粮，结果又把新的责任留在山口。\eventanchor{EM-1445-ESEN-TRIBUTE-CONFLICT}{正统十年朝贡、互市与使团规模的交涉恶化}；使团人数和贡物价值来自交涉中的申报，不能充作中性统计。

\eventanchor{EM-1448-DENG-MAOQI}{正统十三年福建邓茂七起事扩大}，又使县以下的矿役、赋役、迁徙和地方武装进入朝廷的军情语言。“矿徒”“盗贼”是官军分类，不足以容纳参与者的生计背景。中央调兵可以压低一处冲突，却不能证明地方的债务、征役和控制关系已经复原。
\section{皇帝被俘后，谁能发令（1449--1450）}

正统十四年，\eventref{EM-1449-TUMU}{英宗亲征瓦剌、明军在土木堡溃败}。战败、皇帝被俘与主力受损可以确认；路线、军数和责任不能由天顺朝修成的\textit{英宗实录}或后出的个人传记单独裁断。大军的给养、撤退和指挥失去协调，正是多年边警、互市摩擦与中枢信息选择被压进同一次行动的结果。

九月，\eventanchor{EM-1449-BEIJING-DEFENSE}{于谦等组织北京守备并拒绝南迁}。城门、仓储、漕运、京营与民夫使防守成为可能；城墙遗址能显示防御空间，却不能复原每一处兵力和居民选择。守备命令说明朝廷开始调度，不证明粮饷已抵达每一处。与此同时，\eventanchor{EM-1449-JINGTAI-ACCESSION}{郕王朱祁钰即位为景泰帝，尊英宗为太上皇}。这是危机中续接任官、筹饷与守城的正式程序；后世将它预先命名为“夺位”，会抹去当时双重皇权的实际需要。

\eventanchor{ML-1449-009}{也先撤退}后，景泰元年英宗被瓦剌释放，返京后仍以太上皇身份安置。俘虏归来恢复了人身和象征，未恢复直接发令权。明、蒙古叙事均不能提供无歧义的谈判记录；只能说双方都在估量俘虏、边境压力与政治收益，而非替密室中的人补足动机。

\section{复位留下的军政裂缝（1457--1464）}

景泰朝的储位变动使名分紧张延至下一代。天顺元年，\eventanchor{MM-1457-RESTORATION}{石亨、徐有贞、曹吉祥等发动夺门政变，英宗复位}。政变和主要参与者可核，谁先谋划、景泰帝的病况及各人动机却受复辟文书塑形。\eventanchor{MM-1457-YU-QIAN-EXECUTION}{于谦随即以“谋逆”罪被处死}；定罪与处死是事实，罪名首先是清算前朝的司法语言，不能倒写成对其内心和行动的透明证明。

复位也没有让同盟者永久占有宫门与京营。\eventanchor{MM-1460-SHI-HENG-FALL}{天顺四年石亨下狱}，翌年\eventanchor{MM-1461-CAO-QIN-REBELLION}{曹钦在北京兵变、攻东安门后被平定}。地点和结局可靠，参战人数和城市破坏程度多来自战报，不应填成完整图景。两事表明，复位所依靠的武臣、宦官网络一旦进入关键节点，皇帝又必须重划它们的边界。

\eventanchor{MM-1464-XIANZONG-ACCESSION}{天顺八年英宗去世，朱见深即位为宪宗}。储位已经复立，使交接较景泰末年稳定；礼制化的实录仍不能证明宫廷没有摩擦。下一阶段的荆襄、广西、建州与西域，将把这个重新整顿的中枢再次送往山口、边市和地方社会。
